\subsection{20.124: a Rota--Baxter operator in every infinite free rank}
\label{partial:rb}
\problemstatement{20.124}

The question asks for a weight-one Rota--Baxter operator on a
nonabelian free group with image its derived subgroup.
The experiment constructs one in every infinite rank.
The bijective-cocycle framework is prior work
\cite{bardakov-gubarev,guo-lang-sheng}; the construction and
identity can be given directly.

For an infinite cardinal \(\kappa\), partition a free basis
into levels \(X_0,X_1,\ldots\), each of size \(\kappa\).
Let \(F_m\) be free on levels through \(m\).
Both \(F_m\) and \(F_m'\) have cardinality \(\kappa\); for
the lower bound on \(F_m'\), use the distinct commutators
with one fixed basis element. Choose a surjection
\(X_{m+1}\to F_m'\) and define an endomorphism \(\phi\)
by those maps, killing \(X_0\).
It lowers levels, so each element is killed by an iterate.
Moreover \(\phi(F)=\bigcup_m F_m'=F'\).

For any locally terminating endomorphism of any group put
\[
 D(g)=g\phi(g)^{-1},\qquad
 T(g)=g\phi(g)\phi^2(g)\cdots.
\]
Each product is finite for its particular \(g\).
Cancellation gives \(D(T(g))=g=T(D(g))\), so these are inverse
set bijections. Define \(B(g)=\phi(T(g))\), whence
\(B(D(u))=\phi(u)\). For \(g=D(u),h=D(v)\),
\[
 gB(g)hB(g)^{-1}=uv\phi(uv)^{-1}=D(uv).
\]
Therefore
\[
 B(gB(g)hB(g)^{-1})=\phi(uv)=B(g)B(h),
\]
the required identity, and \(\operatorname{im}B=\operatorname{im}\phi=F'\).
If \(\phi(x)=w\), the explicit preimage is \(B(xw^{-1})=w\).
The descendant multiplication is carried from the original
free group by \(D\).
In finite rank an endomorphism image is finitely generated
whereas \(F'\) has infinite rank. This limits this particular
method; it does not rule out another Rota--Baxter map.
See \artifact{research/20.124-infinite-rank-proof.md}.

\subsection{17.118: characteristic exponent subgroups}
\label{partial:17.118}
\problemstatement{17.118}

Let a finite \(p\)-group \(G\) have an exponent-\(p\) subgroup
\(H\) of index \(p\). The question asks for a characteristic
exponent-\(p\) subgroup of index bounded solely in terms of \(p\).
The experiment proves bounds \(p^{c+1}\) if \(H\) has class
at most \(c\), and \(p^c\) if \(G\) has class at most \(c\).
The optimal universal bounds for \(p=2,3\) are \(4,27\).

Here is the class argument. Intersect all automorphic images
of \(H\), obtaining characteristic exponent-\(p\) subgroup
\(K\) with \(G/K\cong\F_p^r\).
Choose \(r\) images \(H_i\) defining independent coordinate
hyperplanes and lifts \(x_i\) of the dual coordinate basis.
Together with generators of \(K\), they generate \(G\), and
omitting \(x_i\) leaves a subgroup inside \(H_i\).
For a class-\(c\) group with this deletion property, \(r>c\)
forces exponent \(p\).
To justify deletion without assuming it is an endomorphism
of \(G\), form the free class-\(c\) group and impose \(w^p=1\)
for every word omitting a distinguished \(x_i\).
The deletion maps descend as commuting idempotents \(e_i\).
Every \(p\)th power lies in \(\bigcap_i\ker e_i\).
On the weight-\(k\) lower-central factor, \(k<r\), the product
\(\prod_i(1-e_i)\) kills each generating bracket, since it
omits some \(x_i\). It is the identity on the common kernel.
Induction thus places that kernel in \(\gamma_r=1\).
This proves the deletion assertion and the bound for \(G\).
If only \(H\) has class \(c\) and \(r>c+1\), every generating
commutator of weight \(c+1\) omits one \(x_i\), hence vanishes
in its \(H_i\). So \(G\) has class at most \(c\), and the same
argument applies.

For \(p=2\), \(H\) is abelian, giving \(4\).
For \(p=3\), Laffey's theorem says that a finite 3-group not
of exponent three has at most a \(7/9\) proportion of cube roots
of one \cite{laffey}.
The union of the independent hyperplanes \(H_i\) has proportion
\(1-(2/3)^r\). At \(r\ge4\) this exceeds \(7/9\).
Thus \(r\le3\), proving \(27\).

For every prime \(p\), a uniform bound must be at least \(p^p\).
Take the associative \(\F_p\)-algebra with basis all words
without repeated letters in \(p\) letters \(x_0,\ldots,x_{p-1}\),
including the empty word, and multiply by concatenation when
the supports are disjoint and by zero otherwise.
Let \(J\) be its positive-degree ideal and \(G=1+J\).
Then \(J^{p+1}=0\). If \(u\) has linear part \(\sum a_ix_i\),
\[
 (1+u)^p=1+\Bigl(\prod_i a_i\Bigr)\Omega,
\]
where \(\Omega\ne0\) is the sum of the distinct length-\(p\)
basis words. The preimage of any coordinate hyperplane
has index \(p\) and exponent \(p\), with class exactly \(p-1\).
The subgroup \(K=1+J^2\) has exponent \(p\) and index \(p^p\).
It is intrinsically characterized by
\[
 K=\{g:(gh)^p=h^p\text{ for every }h\in G\}.
\]
For a nonzero linear part \(a\), choose \(b_j=0\) where
\(a_j\ne0\) and \(b_i\ne-a_i\) for all other coordinates;
then \(\prod b_i=0\ne\prod(a_i+b_i)\), proving the characterization.
Cycling the letters induces a \(p\)-cycle on \(\F_p^p\).
Any nonzero invariant subspace contains a nonzero constant
vector, by taking the last nonzero iterate of \(\sigma-1\).
That vector has a lift with nontrivial \(p\)th power.
The image of a characteristic exponent-\(p\) subgroup must
therefore be zero, forcing it inside \(K\).
This proves sharpness for \(p=2,3\) and the general lower bound.
It does not give unbounded index at a fixed \(p\).
See \artifact{research/17.118-partial-proof.md}; the bounded-class
parameter is not known here to depend only on \(p\ge5\).

\subsection{Other proved restrictions and a lower-bound construction}

\label{partial:21.114}
\problemstatement{21.114}
For \kn{21.114}, write \(a(X)=|X/X'|\); a finite group is
weakly ab-maximal if \(a(H)\le a(X)\) for all \(H\le X\).
The prior general framework is due to Lisi--Sabatini
\cite{lisi-sabatini}. A central-product construction in the
experiment yields order \(8192\) and derived length three,
so a universal derived-length bound, if it exists, is at least three.
Let \(P\) have cyclic central subgroup \(Z\) of order \(p^k\)
with \(Q=P/Z\) weakly ab-maximal. Form the Heisenberg group
\(E\) on \((\Z/p^k\Z)^r\times(\Z/p^k\Z)^r\times\Z/p^k\Z\),
where \(r\ge1\),
with central cross term \(u\cdot v'\), and identify \(E'=Z(E)\)
with \(Z\) in the central product \(G\).
Then \(|G|=|P|p^{2rk}\), \(a(G)=a(Q)p^{2rk}\), and its
derived length is \(\max(2,\operatorname{dl}P)\).
For \(H\le G\), enlarge to \(HZ\), which can only increase \(a(H)\).
In \(H/Z\le Q\times V\), let \(U=(H/Z)\cap V\) and let \(T\)
be the projection on \(Q\). Write
\(|Z/(H'\cap Z)|=p^t\). Then
\[
 a(H)\le p^t|U|a(T)\le p^t|U|a(Q).
\]
The reduction of \(U\) modulo \(p^t\) is isotropic in the
nondegenerate alternating pairing on
\((\Z/p^t\Z)^{2r}\). Character orthogonality gives its size
at most \(p^{rt}\), even when it is not a free module.
Hence \(|U|\le p^{2rk-rt}\), proving \(a(H)\le a(G)\).
Literal finite permutation models supply \(P\) of order \(512\),
derived-series orders \(512,32,2,1\), and central \(Z\) of order
four, with \(Q\) of order \(128\) and \(a(Q)=16\).
Taking \(r=1\) gives the stated group and \(a(G)=256\).
The exact input checks enumerate every subgroup of the two
retained \(Q\)-models; they are recorded in
\artifact{research/21.114-central-products.md}.
That file also proves that central extensions of these particular
outputs still have derived length at most three. Iteration of
this construction does not establish unboundedness.

\label{partial:19.108}
\problemstatement{19.108}
For \kn{19.108}, let \(p\) be odd and \(|P|=p^{10}\), with
abelian normal subgroup \(A\) of order at least \(p^6\).
Then for every irreducible complex \(\chi\) and every
\(\chi(x)\ne0\), one has \(|x|\mid |P|/\chi(1)^2\).
Wilde's prior results cover degree at most \(p^3\), groups
of order at most \(p^9\), and give, for a faithful violating
character, a conjugate of \(\langle x\rangle\) intersecting it
trivially \cite{wilde-supports}.
If \(\chi\) has nontrivial kernel, the order-at-most-\(p^9\)
theorem in the quotient lifts the bound: the order of \(x\)
is at most its quotient order times the kernel order.
Thus a violation requires faithful \(\chi\) of degree \(p^4\).
The bound \(\chi(1)\le[P:A]\) forces \(|A|=p^6\),
and equality in Clifford theory gives \(\chi=\lambda^P\).
Nonvanishing puts \(x\) in the normal subgroup \(A\).
The disjoint conjugate cyclic subgroups then force \(|x|=p^3\)
and \(A=\langle x\rangle\times\langle x\rangle^g
\cong C_{p^3}^2\). The dual orbit of \(\lambda\) has size \(p^4\).
Such an orbit of a \(p\)-subgroup of
\(\GL_2(\Z/p^3\Z)\) is invariant under translation by \(p^2\F_p^2\).
To see this, its image modulo \(p\) has size at most \(p\).
If its image modulo \(p^2\) had size \(p^3\), the congruence
kernel would supply \(I+pM_1,I+pM_2\) with independent
directions at a point. Their \(p\)th powers are
\(I+p^2M_i\) modulo \(p^3\), forcing full \(p^2\)-fibers and
an orbit of size \(p^5\), a contradiction. The image modulo
\(p^2\) thus has size at most \(p^2\); the size \(p^4\)
forces every fiber full. Translation in a coordinate where
\(x\) is primitive multiplies the orbit character sum by a
nontrivial \(p\)th root of unity, so it vanishes, the final
contradiction. Nonsplit extensions are included.
Groups with no such large abelian normal subgroup remain
uncovered; see \artifact{research/19.108-partial.md}.

\label{partial:18.120}
\problemstatement{18.120}
For \kn{18.120}, suppose \(G=AB\), \(A\) abelian and \(B\)
of class at most two. Put \(C=(B')^G=\langle(B')^a:a\in A\rangle\)
and \(K=\langle A,B'\rangle=AC\).
If \(C\) is abelian, \(B_0=K\cap B\) is abelian and \(K=AB_0\).
Indeed an element \(ac\in K\cap B\) centralizes \(B'\);
as \(c\) does too, \(a\) centralizes \(B'\), and commutativity
of \(A\) makes it centralize all its \(A\)-conjugates.
Two such elements commute. The product assertion follows
by writing an element of \(K\) in \(AB\).
If \(G\) has class at most four, conjugates of elements of
\(B'\subseteq\gamma_2(G)\) differ from them by elements
of \(\gamma_3(G)\); all cross brackets have weight at least
five, while \(B'\) is abelian. Thus \(C\) is abelian.
This proves the class-four case without finiteness or
\(A\cap B=1\). The defect-two subnormal case is also a direct
deduction from McCann's earlier lemma \cite{mccann};
see \artifact{research/18.120-partial.md}.

\label{partial:6.47}
\problemstatement{6.47}
For \kn{6.47}, if a binary word defines another group operation
on a class-two group \(G\), collection and the two identity laws
give \(x*y=xy[x,y]^k\). This is associative, with unchanged
inverses and powers and with new commutator
\([x,y]_*=[x,y]^{2k+1}\).
Collect an identity in the free class-two group as a product
of unary powers and binary commutator powers. Setting all but
one, and then all but two, variables equal to one shows that
each collected factor is separately an identity of \(G\).
Unchanged powers and the commutator formula transfer them all
to the new group, which therefore lies in \(\operatorname{var}(G)\).
This is a direct proof of a restricted case within the subject
already studied by Cooper \cite{cooper}, without a novelty claim;
see \artifact{research/6.47-class-two.md}.

\iffullpaper
\subsection{Bounded exclusions with no general conclusion}

The following completed searches are useful negative evidence,
with the stated coverage boundaries.

\begin{itemize}
\item For \kn{20.21}, structural reductions and catalogue
exclusions through order \(1536\) force any counterexample
to have order at least \(3072\). All 26,760 retained kernel
pairs at order \(1536\) were separated from their multiplication
tables; catalogue and subgroup-enumeration coverage remain
dependencies.
\item For \kn{19.20}, all 91,774 nonabelian groups through
order \(511\) gave no equality case and 469 reverse inequalities.
This does not decide the arbitrary-group equality question.
\item For \kn{21.113}, 2,750 ordinary character tables provided
9,850 table/prime cases. Of the corresponding modular cases,
7,573 were available and 2,277 were unavailable. No
counterexample appeared in the available cases.
\item For \kn{21.99}, the retained coverage includes all
4,722 transitive groups of degrees \(2\)--\(20\); direct
word-certificate checks cover the 86 groups of degrees
\(2\)--\(8\). These are different scopes, not interchangeable
verification counts.
\item For \kn{18.43}, exact matrix-trace fingerprints screened
all 3,933,931,043 positive binary necklaces of lengths
\(1\)--\(36\), with no identity candidate.
A future collision would require symbolic certification;
larger lengths remain outside the search.
\end{itemize}

The corresponding reports and scripts are indexed in
\artifact{reports/FINAL_REPORT.md} and \artifact{reports/STATUS.md}.
The absence of a witness in these finite domains is not a
proof for unbounded orders, dimensions, or word lengths.
\fi
